% =====================================================================
% 研究ノート雛形。使い方はリポジトリ直下の AGENTS.md。
% 1プロジェクト1ノート。仮説・識別戦略・結果・判断を、このファイル1つに積み上げていく（正本）。
% 段階4の終わりに tools/lock_pap.py でロックする。ロック前は実データの結果変数を見ない。
% トラックごとの読み替え（README.md「3つのトラック」）:
%   因果推論: 下の章立てのまま使う。
%   構造推定: 「識別戦略と推定」にモデル・推定に使うモーメント・当てはまりを確かめる取り置きデータを書く。
%   推定量・手法: 「データ」をシミュレーションの設定に、「識別戦略と推定」を提案手法と比べる既存手法に読み替える。
% ウェブ版に変換できる書き方は AGENTS.md §3。それ以外のコマンドは portal/build.py が警告を出す。
% =====================================================================
\documentclass[a4paper,10pt]{ltjsarticle}
% 共通プリアンブル。<project>/paper/main.tex から % 共通プリアンブル。<project>/paper/main.tex から % 共通プリアンブル。<project>/paper/main.tex から \input{../../common/lab_preamble}
% 箱は決まったものだけ使う（AGENTS.md §3）。増やすときは AGENTS.md の一覧も直す。
\usepackage{amsmath,amssymb,booktabs,tabularx,longtable,enumitem,xurl,graphicx}
\usepackage[margin=25mm]{geometry}
\usepackage{xcolor}
\usepackage[most]{tcolorbox}
\usepackage{tikz,pgfplots}
\pgfplotsset{compat=1.18}
\usetikzlibrary{arrows.meta,decorations.pathreplacing,positioning,calc}
\usepackage[backend=biber,style=authoryear,natbib=true,maxcitenames=2,doi=true,url=true]{biblatex}
\usepackage[hidelinks]{hyperref}

\definecolor{labBlue}{HTML}{1F5FA8}
\definecolor{labGreen}{HTML}{2E7D4F}
\definecolor{labPurple}{HTML}{6A4C9C}
\definecolor{labOrange}{HTML}{C8641E}
\definecolor{labRed}{HTML}{B03A3A}
\definecolor{labGray}{HTML}{5A6270}

\tcbset{labbase/.style={breakable,enhanced,boxrule=0.6pt,arc=2pt,left=6pt,right=6pt,top=4pt,bottom=4pt,
  fonttitle=\bfseries\small,before skip=8pt,after skip=8pt}}

% 読むための箱
\newtcolorbox{chaptermap}{labbase,colback=labBlue!4,colframe=labBlue,title={この文書の地図}}
\newtcolorbox{keypoint}{labbase,colback=labBlue!4,colframe=labBlue,title={要点}}
\newtcolorbox{term}[1]{labbase,colback=labGray!5,colframe=labGray,title={用語：#1}}
\newtcolorbox{pitfall}{labbase,colback=labRed!4,colframe=labRed,title={落とし穴}}
\newtcolorbox{debate}{labbase,colback=labGray!5,colframe=labGray,title={論争・別の説明}}

% 研究の箱
% 研究仮説：\begin{hypothesis}{H1 短い名前}
\newtcolorbox{hypothesis}[1]{labbase,colback=labGreen!5,colframe=labGreen,title={仮説：#1}}
% 識別戦略：\begin{design}{H1 差の差}
\newtcolorbox{design}[1]{labbase,colback=labBlue!4,colframe=labBlue,title={識別戦略：#1}}
% 推定結果：\begin{result}{H1 の検証（2026-11、確証）}
\newtcolorbox{result}[1]{labbase,colback=labPurple!4,colframe=labPurple,title={結果：#1}}
% 意思決定：続ける／やめる／変える
\newtcolorbox{decision}{labbase,colback=labOrange!6,colframe=labOrange,title={判断（Go/No-Go）}}

% 箱は決まったものだけ使う（AGENTS.md §3）。増やすときは AGENTS.md の一覧も直す。
\usepackage{amsmath,amssymb,booktabs,tabularx,longtable,enumitem,xurl,graphicx}
\usepackage[margin=25mm]{geometry}
\usepackage{xcolor}
\usepackage[most]{tcolorbox}
\usepackage{tikz,pgfplots}
\pgfplotsset{compat=1.18}
\usetikzlibrary{arrows.meta,decorations.pathreplacing,positioning,calc}
\usepackage[backend=biber,style=authoryear,natbib=true,maxcitenames=2,doi=true,url=true]{biblatex}
\usepackage[hidelinks]{hyperref}

\definecolor{labBlue}{HTML}{1F5FA8}
\definecolor{labGreen}{HTML}{2E7D4F}
\definecolor{labPurple}{HTML}{6A4C9C}
\definecolor{labOrange}{HTML}{C8641E}
\definecolor{labRed}{HTML}{B03A3A}
\definecolor{labGray}{HTML}{5A6270}

\tcbset{labbase/.style={breakable,enhanced,boxrule=0.6pt,arc=2pt,left=6pt,right=6pt,top=4pt,bottom=4pt,
  fonttitle=\bfseries\small,before skip=8pt,after skip=8pt}}

% 読むための箱
\newtcolorbox{chaptermap}{labbase,colback=labBlue!4,colframe=labBlue,title={この文書の地図}}
\newtcolorbox{keypoint}{labbase,colback=labBlue!4,colframe=labBlue,title={要点}}
\newtcolorbox{term}[1]{labbase,colback=labGray!5,colframe=labGray,title={用語：#1}}
\newtcolorbox{pitfall}{labbase,colback=labRed!4,colframe=labRed,title={落とし穴}}
\newtcolorbox{debate}{labbase,colback=labGray!5,colframe=labGray,title={論争・別の説明}}

% 研究の箱
% 研究仮説：\begin{hypothesis}{H1 短い名前}
\newtcolorbox{hypothesis}[1]{labbase,colback=labGreen!5,colframe=labGreen,title={仮説：#1}}
% 識別戦略：\begin{design}{H1 差の差}
\newtcolorbox{design}[1]{labbase,colback=labBlue!4,colframe=labBlue,title={識別戦略：#1}}
% 推定結果：\begin{result}{H1 の検証（2026-11、確証）}
\newtcolorbox{result}[1]{labbase,colback=labPurple!4,colframe=labPurple,title={結果：#1}}
% 意思決定：続ける／やめる／変える
\newtcolorbox{decision}{labbase,colback=labOrange!6,colframe=labOrange,title={判断（Go/No-Go）}}

% 箱は決まったものだけ使う（AGENTS.md §3）。増やすときは AGENTS.md の一覧も直す。
\usepackage{amsmath,amssymb,booktabs,tabularx,longtable,enumitem,xurl,graphicx}
\usepackage[margin=25mm]{geometry}
\usepackage{xcolor}
\usepackage[most]{tcolorbox}
\usepackage{tikz,pgfplots}
\pgfplotsset{compat=1.18}
\usetikzlibrary{arrows.meta,decorations.pathreplacing,positioning,calc}
\usepackage[backend=biber,style=authoryear,natbib=true,maxcitenames=2,doi=true,url=true]{biblatex}
\usepackage[hidelinks]{hyperref}

\definecolor{labBlue}{HTML}{1F5FA8}
\definecolor{labGreen}{HTML}{2E7D4F}
\definecolor{labPurple}{HTML}{6A4C9C}
\definecolor{labOrange}{HTML}{C8641E}
\definecolor{labRed}{HTML}{B03A3A}
\definecolor{labGray}{HTML}{5A6270}

\tcbset{labbase/.style={breakable,enhanced,boxrule=0.6pt,arc=2pt,left=6pt,right=6pt,top=4pt,bottom=4pt,
  fonttitle=\bfseries\small,before skip=8pt,after skip=8pt}}

% 読むための箱
\newtcolorbox{chaptermap}{labbase,colback=labBlue!4,colframe=labBlue,title={この文書の地図}}
\newtcolorbox{keypoint}{labbase,colback=labBlue!4,colframe=labBlue,title={要点}}
\newtcolorbox{term}[1]{labbase,colback=labGray!5,colframe=labGray,title={用語：#1}}
\newtcolorbox{pitfall}{labbase,colback=labRed!4,colframe=labRed,title={落とし穴}}
\newtcolorbox{debate}{labbase,colback=labGray!5,colframe=labGray,title={論争・別の説明}}

% 研究の箱
% 研究仮説：\begin{hypothesis}{H1 短い名前}
\newtcolorbox{hypothesis}[1]{labbase,colback=labGreen!5,colframe=labGreen,title={仮説：#1}}
% 識別戦略：\begin{design}{H1 差の差}
\newtcolorbox{design}[1]{labbase,colback=labBlue!4,colframe=labBlue,title={識別戦略：#1}}
% 推定結果：\begin{result}{H1 の検証（2026-11、確証）}
\newtcolorbox{result}[1]{labbase,colback=labPurple!4,colframe=labPurple,title={結果：#1}}
% 意思決定：続ける／やめる／変える
\newtcolorbox{decision}{labbase,colback=labOrange!6,colframe=labOrange,title={判断（Go/No-Go）}}

\addbibresource{../common/refs.bib}

\title{問い：主張の形で書く（例：〇〇政策は△△の□□を高めた）}
\author{担当者名}
\date{2026年11月14日（最終更新）}

\begin{document}
\maketitle

\begin{abstract}
3〜5文で書く。(1) 何を知りたいか、(2) どの仮説をどの識別戦略で確かめるか、(3) 現時点で何が分かったか、(4) 次に何をするか。
推定値には標準誤差か信頼区間と n を付ける。まだ結果がなければ「未検証」と書く。
\end{abstract}

\begin{chaptermap}
\textbf{3行でいうと}
\begin{enumerate}[leftmargin=*]
  \item 背景：なぜ今この問いなのか。
  \item 仮説と識別：何の変動を使って、何の効果を測るか。
  \item 現状：未検証／支持／棄却／保留のどれか。ロックの有無。
\end{enumerate}
\end{chaptermap}

%======================================================================
\section{背景と問い}\label{sec:bg}
この問いについて何が分かっていて、何が分かっていないかを書く。制度の仕組みは、識別に使う部分を特に丁寧に書く。
事実には必ず出典を付ける。

\begin{term}{用語の名前}
読者が初めて見る用語・制度を1つずつ定義する。
\end{term}

%======================================================================
\section{先行研究}\label{sec:prior}
転換点になった論文を5〜10本、表にまとめる。推定値は標準誤差と一緒に写す。「何が未解決か」を必ず書く。
読んだ深さ（本文／要旨／書誌のみ）は \texttt{research/} のメモに残す。

\begin{table}[htbp]
\centering\small
\begin{tabularx}{\linewidth}{llllX}
\toprule
文献 & データ・対象 & 識別戦略 & 推定値（標準誤差） & 何が未解決か \\
\midrule
（著者 年） & 国・期間・n & 差の差など & 0.00 (0.00) & 例：短期の効果のみ \\
\bottomrule
\end{tabularx}
\caption{先行研究の要点。}
\label{tab:prior}
\end{table}

%======================================================================
\section{研究仮説}\label{sec:hyp}
仮説は反証可能な形で書く。推定対象が1文で書けないものは仮説にしない。成功基準と棄却条件は数値で書く。
複数の仮説を検定するなら、主要仮説の指定か多重検定の補正をここで決める。

\begin{hypothesis}{H1 短い名前}
\textbf{仮説}：X は Y を Z 以上変える（符号と大きさ）。\\
\textbf{推定対象}：誰の、どの期間の、どの効果か（例：処置を受けた事業所の、導入後2年間の平均処置効果）。\\
\textbf{識別戦略}：使う変動と手法（詳細は\ref{sec:design}節）。\\
\textbf{データ}：名前・単位・期間・想定 n・秘匿区分。\\
\textbf{根拠}：先行研究の推定値の範囲。理論の予測（筆者の見解と明記）。\\
\textbf{成功基準}：点推定が Z 以上で、95\%信頼区間が 0 を含まない。\\
\textbf{棄却条件}：95\%信頼区間の上限が Z の半分を下回る（「実質的な効果はない」と言える）。\\
\textbf{どちらでもないとき}：保留。検出力不足として、何を足せば決着するかを書く。\\
\textbf{状態}：未検証（未ロック）
\end{hypothesis}

%======================================================================
\section{データ}\label{sec:data}
出所・取得日・版・観測単位・期間・n・欠損・秘匿区分を書く。変数の定義は表にする。
生データの場所と来歴は \texttt{data/README.md} に書き、ここから参照する。
ロック前に見てよいのは、説明変数と標本の構成までである。

%======================================================================
\section{識別戦略と推定}\label{sec:design}
推定式を書き、係数のどれが推定対象に当たるかを示す。

\begin{equation}
  y_{it} = \alpha_i + \lambda_t + \tau D_{it} + \varepsilon_{it}
  \label{eq:main}
\end{equation}

\begin{design}{H1 手法の名前}
\textbf{使う変動}：何が、なぜ外生と考えられるか。\\
\textbf{識別の仮定}：(a) …、(b) …。1つずつ書く。\\
\textbf{脅威と反証テスト}：仮定ごとに「破れていたら何が観察されるか」と、それを見るテストを対応させる。
検定が通っても安全とは限らない。検定の検出力も示す\parencite{roth2022}。\\
\textbf{推論}：標準誤差のクラスタ単位と理由。クラスタ数。多重検定の扱い。\\
\textbf{頑健性}：ロック時点で決めておく代替仕様の一覧。仕様曲線で全体を示す\parencite{simonsohn2020}。
\end{design}

%======================================================================
\section{予想される結果}\label{sec:expected}
結果変数を見る前に、表と図の枠を作る。予想値はすべて「仮定」と書く。
検出力は \texttt{sim/mc\_check.R} で出し、検出できる最小の効果を先行研究の推定値と比べる。
事前に決める意味と限界は \textcite{olken2015} を参照。

\begin{table}[htbp]
\centering\small
\begin{tabular}{lccc}
\toprule
 & (1) 主仕様 & (2) 代替仕様 A & (3) 代替仕様 B \\
\midrule
$\hat\tau$ & （仮定：0.00） & & \\
標準誤差 & & & \\
n／クラスタ数 & & & \\
spec\_id & & & \\
\bottomrule
\end{tabular}
\caption{主な結果の枠（数値は未記入。予想値は筆者の仮定）。}
\label{tab:shell}
\end{table}

%======================================================================
\section{結果}\label{sec:result}
仮説ごとに \texttt{result} 箱を追加する。古い結果は消さずに残し、日付で区別する。
題に「確証」（ロックした仕様）か「探索」（データを見た後の分析）かを必ず書く。

\begin{result}{H1 の検証（未実施）}
\textbf{何をした}：仕様・標本・n・クラスタ。spec\_id。\\
\textbf{何が分かった}：点推定、標準誤差、95\%信頼区間。\\
\textbf{反証テスト}：結果と、そのテストの検出力。\\
\textbf{走らせた仕様の数}：この仮説について \texttt{log/specs.csv} に残る本数。\\
\textbf{判定}：支持／棄却／保留（成功基準・棄却条件に照らして1文で）。
\end{result}

%======================================================================
\section{考察}\label{sec:discussion}
結果が仮説をどう支持・否定したか。別の説明を潰せているか。どこまで一般化できるか。メカニズムは何か。

\begin{pitfall}
有意な結果だけが報告されやすいことは、手法ごとに程度が違う形で確かめられている\parencite{brodeur2020}。
この研究で何本の仕様を試し、どれを本文に載せたかを書く。
\end{pitfall}

%======================================================================
\section{判断と次の一手}\label{sec:decision}
\begin{decision}
\textbf{結論}：Go／No-Go／ピボット。\\
\textbf{理由}：成功基準に照らして1〜2文。\\
\textbf{次の一手}：誰が、いつまでに、何をするか。
\end{decision}

%======================================================================
\section*{更新履歴}
\begin{itemize}
  \item 2026-11-14：作成。
\end{itemize}

\printbibliography[title={参考文献}]
\end{document}
